\documentclass[11pt]{article}
\usepackage[T1]{fontenc}
\usepackage{lmodern}
\usepackage[margin=1in]{geometry}
\usepackage{amsmath,amssymb,amsthm,mathtools}
\usepackage{microtype}
\usepackage{tikz}
\usepackage[colorlinks=true,linkcolor=blue!45!black,citecolor=green!35!black,urlcolor=blue!55!black]{hyperref}
\usepackage[all]{hypcap}
\hypersetup{
  pdftitle={A product counterexample to the simplex maximum for projection-body volume},
  pdfauthor={OpenAI},
  pdfsubject={An elementary counterexample in dimension twenty}
}
\newtheorem{theorem}{Theorem}
\newtheorem{lemma}[theorem]{Lemma}
\newtheorem{proposition}[theorem]{Proposition}
\newtheorem{corollary}[theorem]{Corollary}
\newcommand{\R}{\mathbb R}
\DeclareMathOperator{\conv}{conv}
\DeclareMathOperator{\vol}{vol}
\DeclareMathOperator{\proj}{proj}
\title{A product counterexample to the simplex maximum\\
for projection-body volume}
\author{OpenAI}
\date{September 24, 2026}
\raggedbottom

\begin{document}
\maketitle

\begin{abstract}
The product of two ten-dimensional simplices has larger normalized
projection-body volume than a twenty-dimensional simplex. This gives
a counterexample to Brannen's proposed simplex maximum.
\end{abstract}

\section{Introduction}

Throughout, a convex body in \(\R^d\) is compact and convex with nonempty
interior, and \(d\ge2\). Write \(|K|\) for its \(d\)-dimensional
volume and \(h_C(x)=\max_{y\in C}\langle x,y\rangle\) for the support
function of a compact convex set \(C\). The projection body
\(\Pi_d K\), also written \(\Pi K\) when the dimension is evident, is
defined by
\[
 h_{\Pi K}(u)=\vol_{d-1}\bigl(\proj_{u^\perp}K\bigr),
 \qquad u\in S^{d-1},
\]
extended homogeneously to all of \(\R^d\). Here \(S^{d-1}\) is the
Euclidean unit sphere and \(\proj_{u^\perp}\) is orthogonal projection.
We write \(B_2^j\) for the Euclidean unit ball in \(\R^j\) and
\(\kappa_j=|B_2^j|\).
The functional under consideration is
\begin{equation}\label{eq:functional}
 R_d(K)=\frac{|\Pi K|}{|K|^{d-1}},
 \qquad c_d=\frac{(d+1)d^d}{d!}.
\end{equation}

The extremal values of this affine-invariant functional are classical
questions in convex geometry. Petty established the affine covariance of
the projection-body operator~\cite{Petty}; see also
\cite[Equation~(1)]{Ludwig}. Brannen conjectured in 1996 that simplices
maximize \(R_d\) among all convex bodies~\cite{Brannen}.
The simplex value is \(c_d\), as we verify below. Thus the proposed
simplex upper bound is
\begin{equation}\label{eq:proposed-bound}
 |\Pi K|\le c_d|K|^{d-1}.
\end{equation}
It is the volume of \(\Pi K\) itself that occurs here, not the volume of
its polar: the classical projection inequalities involving the polar
concern a different functional. For \(d\ge3\), Petty's separate minimum
problem asks whether ellipsoids are the only minimizers of the same
functional \(R_d\); see \cite[Section~1]{Saroglou} for these distinctions.

In dimension three, Chen, Feng, Li, Xi, and Xu settle the minimum problem
and also prove
the bound \(R_3(K)\le18=c_3\) for all
three-dimensional convex bodies, with tetrahedra attaining
equality~\cite[Theorems~1.1 and~1.2]{ChenEtAl}. Feng, Hu, Liu, and Xu give
counterexamples to \eqref{eq:proposed-bound} in every dimension
\(d\ge9\), using polytopes with at most \(d+2\)
facets~\cite[Theorem~1.3]{FengEtAl}. Their result already gives a
negative answer to the unrestricted-dimensional simplex-maximum
question. Their construction uses perturbed projections of a cube and
Minkowski's existence theorem~\cite[Section~4]{FengEtAl}. The present
paper supplies a direct Cartesian-product construction and its exact
value: two ten-dimensional simplices suffice.

\begin{theorem}\label{thm:counterexample}
Let \(T_{10}=\conv(0,e_1,\ldots,e_{10})\subset\R^{10}\), where the
\(e_i\) are the coordinate unit vectors, and let
\(K=T_{10}\times T_{10}\subset\R^{20}\). Then
\[
 \frac{R_{20}(K)}{c_{20}}
 =\frac{121\binom{20}{10}}{21\cdot2^{20}}
 =\frac{22\,355\,476}{22\,020\,096}>1.
\]
Thus \(K\) violates \eqref{eq:proposed-bound} in dimension twenty.
\end{theorem}

The exact witness in Theorem~\ref{thm:counterexample} provides an
elementary illustration of the failure of the simplex maximum. The
optimal upper constant and its maximizing bodies remain separate
questions.

The useful mechanism is a product test: for full-dimensional polytopes
\(A\subset\R^r\) and \(B\subset\R^s\), with \(r,s\ge2\),
\[
 R_{r+s}(A\times B)=R_r(A)R_s(B).
\]
Consequently, the product \(c_rc_s\) of the simplex values is a necessary
lower bound for any universal upper constant in dimension \(r+s\). The
dimension-twenty counterexample comes from comparing two copies of the
ten-dimensional simplex value with the twenty-dimensional simplex value.
All identities used in that comparison are proved below.

The growth of upper constants with dimension is a further quantitative
question. For centrally symmetric bodies, Henk gives the lower bound
\(2^d(9/8)^{\lfloor d/3\rfloor}\) for the optimal upper
constant~\cite{Henk}. Iterating our product identity gives an exponential
excess over the simplex benchmark: for some fixed \(\lambda>1\) and
every sufficiently large \(n\), a product of simplices
\(K_n\subset\R^n\) satisfies \(R_n(K_n)\ge\lambda^n c_n\).
Corollary~\ref{cor:exponential-excess} proves this consequence by
repeating ten-dimensional factors and using one additional simplex to
reach each dimension.

Section~\ref{sec:facets} derives the facet formula, the product identity,
and affine covariance. Section~\ref{sec:simplex} computes the simplex
value by expressing its projection body as a cube plus one segment.
Section~\ref{sec:counterexample} combines the identities and gives the
exact integer comparison, then proves the exponential excess in all
sufficiently large dimensions.

\section{Facet vectors and Cartesian products}\label{sec:facets}

We first express projection bodies of polytopes as sums of segments.
The facets of a Cartesian product then separate into two families,
which will make its normalized projection-body volume multiplicative.
The facet formula is the polytopal case of Cauchy's projection formula;
see \cite[Section~1]{Saroglou}. We include its proof
to fix the factor \(1/2\) on which the subsequent constants depend.

We use Minkowski addition \(C+D=\{c+d:c\in C,\ d\in D\}\).
Support functions add under this operation and determine compact convex
sets uniquely. For the latter assertion, if \(x\notin C\) and \(y\in C\)
is nearest to \(x\), minimization along each segment \([y,z]\subset C\)
gives \(\langle x-y,z-y\rangle\le0\) for every \(z\in C\). The linear
functional with vector \(x-y\) therefore separates \(x\) from \(C\).

\begin{samepage}
\begin{lemma}\label{lem:facets}
Let \(P\subset\R^d\), \(d\ge2\), be a full-dimensional convex polytope. For each facet
\(F\), let \(s_F\) be its \((d-1)\)-dimensional volume and \(\nu_F\) its
outward unit normal. Then
\begin{equation}\label{eq:facet}
 h_{\Pi P}(u)=\frac12\sum_Fs_F|\langle\nu_F,u\rangle|,
 \qquad u\in\R^d,
\end{equation}
and therefore
\begin{equation}\label{eq:zonotope}
 \Pi P=\sum_F
 \left[-\frac{s_F\nu_F}{2},\frac{s_F\nu_F}{2}\right].
\end{equation}
\end{lemma}
\end{samepage}

\begin{proof}
First take \(u\) to be a unit vector. Orthogonal projection from the
hyperplane of \(F\) to \(u^\perp\) has absolute Jacobian
\(|\langle\nu_F,u\rangle|\). Indeed, for an orthonormal tangent basis
\(v_1,\ldots,v_{d-1}\), that Jacobian is
\(|\det(v_1,\ldots,v_{d-1},u)|\), which equals the stated inner product
in absolute value.

Apart from a set of \((d-1)\)-dimensional measure zero, a point in the
projection of \(P\) lies on the projections of exactly two facet
interiors: the entry and exit points of its line parallel to \(u\).
To see that the exceptions are null, discard the boundary of the
projection and the projections of all faces of dimension at most
\(d-2\). Also discard projections of facets whose hyperplanes are parallel to
\(u\); these projections have dimension at most \(d-2\).
Every remaining line meets \(P\) in a nondegenerate interval with its
two endpoints in facet interiors. Integrating this multiplicity gives
\eqref{eq:facet} for unit \(u\), and homogeneity gives the general case.

The support function of \([-v/2,v/2]\) is
\(\tfrac12|\langle v,u\rangle|\). Adding these support functions proves
\eqref{eq:zonotope}.
\end{proof}

\begin{proposition}\label{prop:product}
Let \(A\subset\R^r\) and \(B\subset\R^s\) be full-dimensional convex
polytopes, where \(r,s\ge2\). In the orthogonal product space
\(\R^{r+s}=\R^r\times\R^s\),
\begin{equation}\label{eq:product-body}
 \Pi(A\times B)=(|B|\Pi A)\times(|A|\Pi B).
\end{equation}
In particular,
\begin{equation}\label{eq:multiplicative}
 R_{r+s}(A\times B)=R_r(A)R_s(B).
\end{equation}
\end{proposition}

\begin{proof}
The face exposed by a linear functional \((u,v)\) on \(A\times B\)
is the product of the faces exposed by \(u\) and \(v\).
Dimensions add under Cartesian products. A nonzero linear functional
exposes a proper face of a full-dimensional body, so if both \(u\) and
\(v\) are nonzero, the product face has codimension at least two.
A facet therefore has one facet factor and one whole factor.
Thus the facets are \(F\times B\)
and \(A\times G\), with area-normal vectors
\[
 (|B|s_F\nu_F,0)
 \quad\hbox{and}\quad
 (0,|A|s_G\nu_G),
\]
respectively. The area factors are products because the two coordinate
spaces are orthogonal. Lemma~\ref{lem:facets} now gives
\[
 h_{\Pi(A\times B)}(u,v)
 =|B|h_{\Pi A}(u)+|A|h_{\Pi B}(v).
\]
This is the support function of the right side of
\eqref{eq:product-body}, proving that identity. Hence
\[
 \begin{split}
 R_{r+s}(A\times B)
 &=\frac{|B|^r|A|^s|\Pi A||\Pi B|}
         {(|A||B|)^{r+s-1}}\\
 &=\frac{|\Pi A|}{|A|^{r-1}}
   \frac{|\Pi B|}{|B|^{s-1}},
 \end{split}
\]
as required.
\end{proof}

The product identity reduces our comparison to the values of its factors.
To use simplices as a dimension-dependent benchmark, we also need to
know that all full-dimensional simplices of a fixed dimension have the
same value. The following affine covariance formula supplies that fact.

\begin{lemma}\label{lem:affine}
Let \(P\subset\R^d\), \(d\ge2\), be a convex body. For every
invertible linear map \(L\colon\R^d\to\R^d\),
\[
 \Pi(LP)=|\det L|L^{-t}\Pi P,\qquad R_d(LP)=R_d(P),
\]
where \(L^{-t}\) denotes the inverse transpose. Translations also
preserve \(R_d\).
\end{lemma}

\begin{proof}
First suppose that \(P\) is a polytope.
The transformed outward unit normal of \(F\) is
\(L^{-t}\nu_F/\|L^{-t}\nu_F\|\). Its area is
\[
 s_{LF}=|\det L|\,\|L^{-t}\nu_F\|\,s_F.
\]
For the area formula, take a prism of height one over \(F\), extending
in direction \(\nu_F\). Its image has volume \(|\det L|s_F\) and height
\(1/\|L^{-t}\nu_F\|\) over \(LF\); dividing volume by height gives the
formula. Thus each area-normal vector transforms by
\(|\det L|L^{-t}\), including when \(\det L<0\).
Apply \eqref{eq:zonotope}. The absolute determinant of this map is
\(|\det L|^{d-1}\), exactly the scaling of \(|LP|^{d-1}\).
Translation invariance follows directly from projection volumes.

For completeness, we justify the passage to an arbitrary convex body
and the existence of its projection body. Translate an interior point
of \(P\) to the origin, choose \(R>0\) with \(P\subset RB_2^d\), and
take increasing inscribed polytopes \(P_m\) containing a fixed ball
\(aB_2^d\), \(a>0\), such that
\[
 P_m\subset P\subset P_m+\varepsilon_mB_2^d
 \subset(1+\varepsilon_m/a)P_m,
 \qquad \varepsilon_m\longrightarrow0.
\]
Take, for example, convex hulls of cumulative finite nets in \(P\) and
a fixed inscribed polytope containing \(aB_2^d\). Put
\(\lambda_m=1+\varepsilon_m/a\), and for unit \(u\) write
\(F(u)=\vol_{d-1}(\proj_{u^\perp}P)\) and
\(f_m(u)=h_{\Pi P_m}(u)\). Projection and volume monotonicity give
\[
 0\le F(u)-f_m(u)
 \le (\lambda_m^{d-1}-1)\kappa_{d-1}R^{d-1}.
\]
Thus \(f_m\to F\) uniformly. The increasing bodies \(\Pi P_m\) lie
in a fixed ball and contain \(\kappa_{d-1}a^{d-1}B_2^d\), so their
union closure is a convex body with support function \(F\), namely
\(\Pi P\). The support inequalities also give
\[
 \Pi P_m\subset\Pi P\subset\lambda_m^{d-1}\Pi P_m,
 \qquad
 |\Pi P_m|\le|\Pi P|\le\lambda_m^{d(d-1)}|\Pi P_m|.
\]
Likewise \(|P_m|\le|P|\le\lambda_m^d|P_m|\). Applying \(L\) to
the original sandwich gives the same convergence for \(LP_m\).
Passing to the limit in the polytope identity proves covariance and,
by these volume inequalities, the asserted invariance of \(R_d\).
\end{proof}

\section{The simplex value}\label{sec:simplex}

We now compute the value to which the product will be compared.
Only one extra segment beyond a coordinate cube is needed to describe
the projection body of a simplex.

\begin{proposition}\label{prop:simplex}
For \(d\ge2\), the standard simplex
\(T_d=\conv(0,e_1,\ldots,e_d)\subset\R^d\) satisfies
\[
 |T_d|=\frac1{d!},\qquad
 |\Pi T_d|=\frac{d+1}{((d-1)!)^d},\qquad
 R_d(T_d)=c_d.
\]
Every full-dimensional \(d\)-simplex has the same value \(c_d\).
\end{proposition}

\begin{proof}
The simplex volume follows by iterated base and height. Put
\(w=e_1+\cdots+e_d\). The coordinate facets have area-normal vectors
\(-e_i/(d-1)!\). The remaining facet has unit normal \(w/\sqrt d\).
Dropping its last coordinate projects it onto the standard
\((d-1)\)-simplex with Jacobian \(1/\sqrt d\); its area-normal vector is
therefore \(w/(d-1)!\). In the centered segments of
Lemma~\ref{lem:facets}, changing the sign of a generator changes no
segment. Translating each segment to start at the origin therefore shows
that \(\Pi T_d\) is a translate of
\[
 \frac{[0,1]^d+[0,w]}{(d-1)!}.
\]

To compute the numerator's volume, let \(Q=[0,1]^d\).
Every nonempty fiber of \(Q\) parallel to \(w\) is an interval, and
adding \([0,w]\) increases its length by \(\|w\|\).
The base projection is unchanged. Fubini therefore gives
\[
 |Q+[0,w]|
 =|Q|+\|w\|\vol_{d-1}(\proj_{w^\perp}Q)
 =1+h_{\Pi Q}(w).
\]
The cube has two facets of area one in each coordinate direction,
so \eqref{eq:facet} gives \(h_{\Pi Q}(w)=d\).
Thus \(|Q+[0,w]|=d+1\); Figure~\ref{fig:extrusion} illustrates the
two-dimensional case. Dilation by \(1/(d-1)!\) and division by
\(|T_d|^{d-1}\) now yield
\[
 R_d(T_d)=\frac{(d+1)(d!)^{d-1}}{((d-1)!)^d}
         =\frac{(d+1)d^d}{d!}.
\]
Every full-dimensional simplex is an invertible affine image of
\(T_d\), so Lemma~\ref{lem:affine} proves the final assertion.
\end{proof}

\begin{figure}[tb]
\centering
\begin{tikzpicture}[scale=1.5, line join=round]
 \fill[black!7] (0,0) rectangle (1,1);
 \fill[blue!13] (1,0)--(2,1)--(2,2)--(1,1)--cycle;
 \fill[yellow!23] (0,1)--(1,1)--(2,2)--(1,2)--cycle;
 \draw[thick] (0,0)--(1,0)--(2,1)--(2,2)--(1,2)--(0,1)--cycle;
 \draw[thin] (0,1)--(1,1)--(1,0);
 \draw[thin] (1,1)--(2,2);
 \node at (.5,.5) {\(Q\)};
 \node at (1.5,1) {\(1\)};
 \node at (1,1.5) {\(1\)};
 \node[below left] at (0,0) {\(0\)};
 \node[above right] at (2,2) {\((2,2)\)};
 \draw[->,thick] (2.6,.3)--(3.6,1.3);
 \node[right] at (3.15,.65) {\(w=(1,1)\)};
\end{tikzpicture}
\caption{The planar extrusion \(Q+[0,(1,1)]\), with
\(Q=[0,1]^2\), consists of the unit square and two parallelograms of
area one. Its area is \(3\). The fiber argument in
Proposition~\ref{prop:simplex} proves the corresponding volume \(d+1\)
in every dimension.}
\label{fig:extrusion}
\end{figure}

\section{The counterexample and exponential excess}\label{sec:counterexample}

The geometric work is complete. The product formula turns the comparison
in dimension twenty into a calculation with two simplex constants.

\begin{proof}[Proof of Theorem~\ref{thm:counterexample}]
The product \(K=T_{10}\times T_{10}\) is compact and convex with
nonempty interior in \(\R^{20}\). Propositions~\ref{prop:product}
and~\ref{prop:simplex} give
\[
 R_{20}(K)=R_{10}(T_{10})^2
         =\left(\frac{11\cdot10^{10}}{10!}\right)^2.
\]
Consequently,
\[
 \frac{R_{20}(K)}{c_{20}}
 =\frac{121\cdot10^{20}\cdot20!}
        {(10!)^2\cdot21\cdot20^{20}}
 =\frac{121\binom{20}{10}}{21\cdot2^{20}}
 =\frac{22\,355\,476}{22\,020\,096}>1.
\]
Here \(\binom{20}{10}=184\,756\), and the exact integer certificate is
\[
 121\cdot184\,756-21\cdot1\,048\,576=335\,380>0.
\]
This is an exact strict violation of \eqref{eq:proposed-bound}.
\end{proof}

More generally, for \(r,s\ge2\), the same two propositions give the
explicit product test
\[
 \frac{R_{r+s}(T_r\times T_s)}{c_{r+s}}
 =\frac{(r+1)(s+1)}{r+s+1}
   \binom{r+s}{r}\frac{r^rs^s}{(r+s)^{r+s}}.
\]
Thus any universal upper constant in dimension \(r+s\) must be at least
\(c_rc_s\).

Iterating the product test gives the following quantitative consequence.

\begin{corollary}\label{cor:exponential-excess}
There exist a constant \(\lambda>1\) and an integer \(N\ge20\) such that, for every
integer \(n\ge N\), some full-dimensional convex polytope
\(K_n\subset\R^n\) satisfies
\[
 \frac{R_n(K_n)}{R_n(T_n)}\ge\lambda^n.
\]
The polytopes \(K_n\) may be chosen as Cartesian products of simplices.
\end{corollary}

\begin{proof}
We use repeated ten-dimensional factors to obtain exponential growth,
and one factor of dimension between ten and nineteen to reach every
sufficiently large dimension. For \(r\in\{0,\ldots,9\}\), put
\(d_r=10+r\). Every integer \(n\ge20\) has a unique representation
\(n=10k+d_r\) with \(k\ge1\). In the orthogonal coordinate product
space, take
\[
 K_n=T_{10}^{\,k}\times T_{d_r}.
\]
Here \(T_{10}^{\,k}\) denotes the \(k\)-fold Cartesian product. All
factors and their successive products are full-dimensional convex
polytopes, with each factor dimension at least ten. Thus
Propositions~\ref{prop:product} and~\ref{prop:simplex} apply iteratively
and give
\[
 R_n(K_n)=c_{10}^k c_{d_r}.
\]

To compare this value with \(c_n\), we bound the simplex constants by
\(n+1\) times a fixed exponential, whose rate is smaller than that
supplied by the ten-dimensional factors. Put \(\beta=11/4\).
The exponential series gives
\[
 e=\frac52+\sum_{j=3}^{\infty}\frac1{j!}
 <\frac52+\frac16\sum_{j=0}^{\infty}3^{-j}
 =\beta,
\]
since the successive ratios in the tail are strictly less than \(1/3\).
The positive term \(n^n/n!\) in the series for \(e^n\) is strictly
smaller than the full sum, so
\[
 c_n<(n+1)e^n<(n+1)\beta^n\qquad(n\ge2).
\]
For the fixed ten-dimensional block, exact arithmetic yields
\[
 c_{10}=\frac{17\,187\,500}{567},\qquad
 A:=\frac{c_{10}}{\beta^{10}}
 =\frac{1\,638\,400\,000\,000}{1\,336\,956\,340\,797}>1.
\]
Choose a fixed \(\lambda\) with \(1<\lambda<A^{1/10}\). Combining
these estimates with the formula for \(R_n(K_n)\) gives
\[
 \frac{R_n(K_n)}{R_n(T_n)\lambda^n}
 >c_{d_r}\beta^{-d_r}\lambda^{-d_r}
   \frac{(A/\lambda^{10})^k}{10k+d_r+1}.
\]
For each fixed \(r\), the right side tends to infinity as \(k\to\infty\),
because \(A/\lambda^{10}>1\). A threshold can therefore be chosen in
each of the ten residue classes; their maximum, enlarged to be at least
twenty, gives one \(N\) for all \(n\ge N\).
\end{proof}

\begin{thebibliography}{7}

\bibitem{Brannen}
N.~S.~Brannen,
\emph{Volumes of projection bodies},
Mathematika \textbf{43} (1996), no.~2, 255--264.
\href{https://doi.org/10.1112/S002557930001175X}{doi:10.1112/S002557930001175X}.

\bibitem{ChenEtAl}
S.~Chen, Y.~Feng, Y.~Li, D.~Xi, and L.~Xu,
\emph{Petty's conjectured projection inequality in dimension three},
preprint, 2026, MathSciDoc:2608.23001.
\href{https://archive.ymsc.tsinghua.edu.cn/pacm_download/741/12780-pettyinequality.pdf}{Full text}.

\bibitem{FengEtAl}
Y.~Feng, S.~Hu, W.~Liu, and L.~Xu,
\emph{On the reverse projection inequality},
preprint, 2026, MathSciDoc:2608.23002.
\href{https://archive.ymsc.tsinghua.edu.cn/pacm_download/743/12781--2026.8.26.pdf}{Full text}.

\bibitem{Henk}
M.~Henk,
\emph{Note on projection bodies of zonotopes with \(n+1\) generators},
Acta Mathematica Scientia \textbf{45} (2025), 96--103.
\href{https://doi.org/10.1007/s10473-025-0107-9}{doi:10.1007/s10473-025-0107-9}.

\bibitem{Ludwig}
M.~Ludwig,
\emph{Projection bodies and valuations},
Advances in Mathematics \textbf{172} (2002), no.~2, 158--168.
\href{https://doi.org/10.1016/S0001-8708(02)00021-X}{doi:10.1016/S0001-8708(02)00021-X}.

\bibitem{Petty}
C.~M.~Petty,
\emph{Projection bodies},
in \emph{Proceedings of the Colloquium on Convexity (Copenhagen, 1965)},
K\o benhavns Universitet Matematiske Institut, Copenhagen, 1967,
pp.~234--241.

\bibitem{Saroglou}
C.~Saroglou,
\emph{On the shape of a convex body with respect to its second projection body},
Advances in Applied Mathematics \textbf{67} (2015), 55--74.
\href{https://doi.org/10.1016/j.aam.2015.03.004}{doi:10.1016/j.aam.2015.03.004}.
Preprint:
\href{https://arxiv.org/abs/1409.4347v2}{arXiv:1409.4347v2}.

\end{thebibliography}

\end{document}
